\section{Design and Implementation}
\eva{} provides a unified assessment architecture that decouples workload definition from execution infrastructure, supports heterogeneous model providers and agent frameworks through a standardized API layer, and records both outcome and process metrics as first-class analytical objects.

\subsection{Assessment Framework}

\eva{} consists of five major components (Figure~\ref{fig:architecture}):

\textbf{Task Orchestrator.}
Responsible for workload scheduling, dependency resolution, resurrection triggering, and evaluation pipeline management.
The orchestrator maintains the serial task chain state, monitors per-task pass/fail status, and triggers resurrection when a task falls below the 60\% passing threshold.

\textbf{Agent Runtime Layer.}
Supports multiple agent frameworks through a unified interface, each representing a distinct interaction paradigm:

OpenHands SDK: Full agent framework with unlimited self-loop iteration, skill injection, and MCP tool integration.
Claude Code: File I/O and command execution agent with multi-turn conversation.
Codex CLI: Autonomous coding agent with sandboxed execution.
Kimi CLI: Conversational coding agent with plan management and background task execution.

All backends share a common task context format and produce standardized result structures, enabling direct cross-backend comparison.

\textbf{Model Access Layer (\textsc{AIhub}).}
A unified API gateway (\url{https://aihub.arcsysu.cn}) providing standardized OpenAI-compatible access to 21 large language models spanning proprietary and open-weight families.
This design eliminates API-specific integration overhead and ensures that evaluation differences reflect model and scaffold capabilities rather than infrastructure variability.

\textbf{Execution Environment.}
Each evaluation run operates in an isolated workspace containing a complete copy of the target repository with all dependencies (ANTLR, LLVM 14, pybind11) pre-installed.
For \zzc{\Progressive}{}, each run additionally launches a fresh container environment that is destroyed after completion, ensuring zero cross-run contamination.
The workspace is configured via CMake with a fixed toolchain; scoring scripts are deterministic.

\textbf{Observation and Leaderboard Layer.}
Collects and aggregates \warn{evaluation} results, computes utility-oriented metrics, and maintains dynamic leaderboards.
All raw evaluation data (scores, logs, process metrics) is preserved for reproducibility.

\begin{figure}[t]
\centering
    \includegraphics[width=\linewidth]{figures/fig1_sys_arch.png}
\caption{The \eva{} assessment framework architecture. Five layers from top to bottom: Observation \& Leaderboard, Task Orchestrator, Agent Runtime, Model Access (AIhub), and Execution Environment.\warn{最下层把YATCC和YATCC-Hard分开写了，但是这两个按我理解应该是一个执行环境？}}
\label{fig:architecture}
\end{figure}

\subsection{Assessment Tasks}

\eva{} grounds its \warn{evaluation} in compiler construction, a domain chosen for its strict causal dependencies, well-defined intermediate artifacts, and objective automated scoring.
Compiler engineering provides an ideal \warn{evaluation} substrate because: (1)~the pipeline has \emph{strict causal dependencies}---each stage consumes the output of its predecessor; (2)~intermediate artifacts are \emph{well-defined and verifiable}---token streams, parse trees, IR modules, and assembly files have clear structural specifications; (3)~the project includes \emph{automated scoring scripts} with comprehensive test suites.

\subsubsection{Workloads: Compiler Construction as Realistic \warn{Evaluation}}
We provide \zzc{a representative workload}{two complementary workloads spanning a difficulty spectrum}:

\zzc{
{\bf \workload{}: Scaffolded Implementation}.
The original \Incremental{} workload~\cite{yatcc2024} organizes six tasks along the compiler construction pipeline (Table~\ref{tab:tasks}).
Each task provides \emph{scaffolded starter code}---the agent fills in missing implementation details within a provided framework.
Crucially, \workload{} persists the full execution context and runtime environment across all tasks, simulating the real-world agentic workflow where agents operate within a single, unified environment to complete tasks sequentially. This design eliminates artificial context resets between tasks, providing a more realistic evaluation of continuous system development capabilities.
This workload measures an agent's ability to \emph{complete} a partially-implemented system incrementally.
}{}

% {\bf \zzc{\workload{}}{\yatcc{}-Hard}: Full From-Scratch Implementation}.
% \zzc{\workload{}}{\yatcc{}-Hard} removes all code logic from the original tasks, retaining only file dependency relationships to ensure automated scripts remain functional.
% Agents receive \emph{no compiler construction knowledge} and \emph{no implementation framework}---they must implement the full compiler from scratch.
% The task descriptions are reduced to experiment requirements and scoring script usage instructions.
% This workload measures an agent's ability to \emph{construct} a complete system from minimal specification.




\begin{table*}[!t]
\centering
\caption{\eva{} task chain. Each task consumes the output artifact of its predecessor, forming a strict serial dependency.}
\small
\begin{tabular}{clll}
\toprule
\textbf{Task} & \textbf{Name} & \textbf{Objective} & \textbf{Input Artifact} \\
\midrule
0 & Environment Setup & Build \& verify toolchain & --- \\
1 & Lexer & Tokenize C source code & Preprocessed source \\
2 & Parser & Construct AST/ASG & Token stream (from T1) \\
3 & IR Generation & Emit LLVM IR & AST/JSON (from T2) \\
4 & IR Optimization & Implement LLVM Pass & Raw IR (from T3) \\
5 & \zzc{Assembly}{Code} Generation & Produce RV64 assembly & Optimized IR (from T4) \\
\bottomrule
\end{tabular}
\label{tab:tasks}
\end{table*}

% {\bf \zzc{\Incremental{}}{\yatcc{}}: Scaffolded Implementation}.
% The original \zzc{\Incremental{}}{\yatcc{}} workload~\cite{yatcc2024} organizes six tasks along the compiler construction pipeline (Table~\ref{tab:tasks}).
% Each task provides \emph{scaffolded starter code}---the agent fills in missing implementation details within a provided framework.
% This workload measures an agent's ability to \emph{complete} a partially-implemented system.

% The difficulty gap between \zzc{\Progressive{}}{\yatcc{}} and \zzc{\Incremental{}}{\yatcc{}-Hard} quantifies the value of scaffolding in agentic coding tasks and reveals which models depend on provided structure versus which can operate autonomously.

\begin{figure*}[!t]
\centering
\includegraphics[width=0.7\textwidth]{figures/fig2_task_rec.png}
\caption{Serial dependency chain with resurrection mechanism. Six tasks (T0--T5) form a strict pipeline where each task's output becomes the next task's input. When a task fails (score $<$ 60\%), the resurrection mechanism injects the golden intermediate artifact, allowing downstream evaluation to continue.}
\label{fig:task_chain}
\end{figure*}

\subsubsection{Resurrection Protocol}

The \resurrection{} mechanism addresses the fundamental evaluation challenge of serial workloads: when an agent fails at Task~$k$, downstream tasks become unreachable, and the benchmark can only report ``failure at $k$'' without diagnosing whether the agent could succeed at Tasks~$k+1, \ldots, 5$ given correct prerequisites.

\warn{\textbf{Definition.}}
When Task~$k$ scores below the passing threshold (60\%), \resurrection{} is triggered:
(1)~the agent's failed artifact is replaced with the \emph{golden intermediate artifact}---the correct output produced by the reference implementation;
(2)~the build system is reconfigured via \texttt{config.cmake} (\texttt{TASK\_k\_REVIVE=ON});
(3)~the agent continues to Task~$k+1$ with the corrected state.

\textbf{Design rationale.}
Resurrection is an \emph{\warn{evaluation} protocol}, not a scoring adjustment.
It decomposes the compound failure ``cannot reach Task~$k+1$'' into two independent measurements:
(1)~\emph{predecessor failure}: the agent could not produce a passing artifact at Task~$k$;
(2)~\emph{downstream capability}: given correct prerequisites, can the agent succeed at Task~$k+1$?

This decomposition enables node-level pass rate reporting even when the full pipeline is broken, significantly improving diagnostic resolution.
Agents are \emph{not informed} about resurrection---the mechanism is triggered externally by the \zzc{task}{evaluation} orchestrator, eliminating moral hazard.

\subsection{Utility-Oriented Multi-Dimensional Metrics}

\eva{} proposes a multi-dimensional assessment metric system that moves beyond simple accuracy scores toward utility-oriented measurement.
Assessment metrics should measure not just \emph{whether} an agent succeeds, but \emph{how efficiently} it succeeds and \emph{how} it fails.

\zzc{To capture these complementary dimensions of agent utility, we define two foundational scoring paradigms that govern how task-level scores are computed, upon which all higher-level aggregate metrics are built.}{}

\zzc{\textbf{End-to-End Production Readiness Scoring (E2E-PRS):}
Simulates real-world production requirements where partial functionality is unacceptable:
\begin{itemize}
    \item Tasks 0–3 and 5 are binary pass/fail (100/0 points)
    \item Any task failure cascades to 0 points for all subsequent tasks
    \item Task 4 (Performance) requires perfect functional verification first; only then is performance score awarded
\end{itemize}
}{}


\zzc{\textbf{Granular Capability Profiling Scoring (GCPS):}
Measures incremental progress and identifies specific strengths/weaknesses:
\begin{itemize}
    \item All tasks receive partial credit (0–100 points) based on requirement completion
    \item No cascading failures; subsequent tasks are evaluated independently
    \item Task 4 still requires functional verification for performance scoring
\end{itemize}
}{}


\textbf{Outcome metrics.}
We define three primary outcome metrics.
\emph{Mean Reward (MR)} is the weighted average of task scores with resurrection bonus: $\text{MR} = \sum s_i w_i b_i / \sum w_i b_i$, where $w_i = [0.05, 0.20, 0.20, 0.15, 0.30, 0.10]$ and $b_i = 1.2$ (no resurrection) or $1.0$ (resurrected).
\emph{Pass Score (PS)} is the weighted pass rate: $\text{PS} = \sum p_i b_i^{\text{pass}} / (6 \times 1.5) \times 100$, where $p_i = 1$ if $s_i \geq 60$ and $b_i^{\text{pass}} = 1.5$ (no resurrection) or $1.0$.
\emph{Pipeline Score (PL)} is the simple average: $\text{PL} = \frac{1}{6}\sum s_i$.

\textbf{Process metrics.}
Beyond outcome scores, we record token consumption, interaction turns, command executions, retry counts, and elapsed wall-clock time per task and total.

\textbf{Failure taxonomy.}
We categorize agent failures into five types:
(1)~\emph{predecessor failure}---downstream unreachable due to missing artifact;
(2)~\emph{task-local reasoning failure}---cannot solve despite correct prerequisites;
(3)~\emph{execution failure}---build/dependency errors;
(4)~\emph{process collapse}---max-turn exhaustion, search loops, repetitive debugging without progress;
(5)~\emph{cost overrun}---budget exhausted.

This taxonomy enables systematic analysis of \emph{how} agents fail, not just \emph{whether} they fail.
For instance, process collapse is more prevalent in framework-based agents with unlimited iteration budgets, suggesting that scaffold design affects failure mode distribution.